\begin{longtable}{@{}p{6.1cm}rrrr@{}}
\caption{Chaque variante contre le Top quadratique, sur les mêmes coureurs, en points : la part des duels bien prédits de 2012 à 2024 sur toutes les courses, sur les courses nationales, puis depuis juillet 2025 ; la part des écarts de temps expliquée (R$^2$) de 2012 à 2024. Un astérisque marque un écart significatif : son intervalle à 95 \% ne contient pas 0.}\label{tab:variants}\\
\toprule Variante & Duels & Nationales & Depuis 2025 & R$^2$ \\ \midrule \endfirsthead
\toprule Variante & Duels & Nationales & Depuis 2025 & R$^2$ \\ \midrule \endhead
\multicolumn{5}{@{}l}{\textsc{Forêt}} \\
\multicolumn{5}{@{}l}{\emph{Références}} \\
CN officiel de la FFCO & $-$1,78\,\textbf{*} & $-$1,18\,\textbf{*} & $-$1,19\,\textbf{*} & $-$9,48\,\textbf{*} \\
Top linéaire : règle FFCO pour le circuit, CN en moyenne ordinaire & $-$0,11\,\textbf{*} & $-$0,05\,\textbf{*} & $-$0,08\,\textbf{*} & $-$1,32\,\textbf{*} \\
\multicolumn{5}{@{}l}{\emph{Exposant sur le temps, règle FFCO pour le circuit}} \\
exposant 1,5 & +0,13\,\textbf{*} & +0,02 & +0,09\,\textbf{*} & +0,75\,\textbf{*} \\
exposant 2 & +0,20\,\textbf{*} & +0,00 & +0,19\,\textbf{*} & +1,55\,\textbf{*} \\
exposant 3 & +0,07\,\textbf{*} & $-$0,21\,\textbf{*} & $-$0,09 & +1,81\,\textbf{*} \\
exposant 2 pour le circuit seulement & +0,15\,\textbf{*} & $-$0,01 & +0,14\,\textbf{*} & +0,44\,\textbf{*} \\
\multicolumn{5}{@{}l}{\emph{Règle FFCO pour le circuit, autre moyenne pour le CN}} \\
CN en moyenne cubique & +0,08\,\textbf{*} & +0,03\,\textbf{*} & +0,07\,\textbf{*} & +0,61\,\textbf{*} \\
\multicolumn{5}{@{}l}{\emph{Autres mesures du circuit}} \\
géométrique, deux tiers, CN en moyenne ordinaire & $-$0,39\,\textbf{*} & $-$0,18\,\textbf{*} & $-$0,43\,\textbf{*} & $-$2,91\,\textbf{*} \\
géométrique, deux tiers, CN en moyenne quadratique & $-$0,25\,\textbf{*} & $-$0,09\,\textbf{*} & $-$0,26\,\textbf{*} & $-$1,47\,\textbf{*} \\
arithmétique, tout le peloton, CN en moyenne ordinaire & +0,33\,\textbf{*} & +0,00 & +0,22\,\textbf{*} & +1,32\,\textbf{*} \\
arithmétique, tout le peloton, CN en moyenne quadratique & +0,32\,\textbf{*} & $-$0,04 & +0,21\,\textbf{*} & +2,15\,\textbf{*} \\
arithmétique, tout le peloton, CN en moyenne cubique & +0,27\,\textbf{*} & $-$0,12\,\textbf{*} & +0,17\,\textbf{*} & +2,13\,\textbf{*} \\
arithmétique, 85 \% les plus rapides, CN en moyenne quadratique & +0,39\,\textbf{*} & +0,12\,\textbf{*} & +0,25\,\textbf{*} & +2,05\,\textbf{*} \\
géométrique, 85 \% les plus rapides, CN en moyenne cubique & +0,27\,\textbf{*} & +0,10\,\textbf{*} & +0,25\,\textbf{*} & +1,59\,\textbf{*} \\
géométrique atténuant les temps aberrants (Huber), quadratique & +0,23\,\textbf{*} & +0,07\,\textbf{*} & +0,14\,\textbf{*} & +1,35\,\textbf{*} \\
\multicolumn{5}{@{}l}{\emph{Moyenne géométrique sur tout le peloton}} \\
CN en moyenne ordinaire & +0,27\,\textbf{*} & +0,09\,\textbf{*} & +0,17\,\textbf{*} & +0,72\,\textbf{*} \\
CN en moyenne quadratique & +0,35\,\textbf{*} & +0,13\,\textbf{*} & +0,27\,\textbf{*} & +1,94\,\textbf{*} \\
CN en moyenne cubique (retenu en forêt) & +0,40\,\textbf{*} & +0,12\,\textbf{*} & +0,31\,\textbf{*} & +2,37\,\textbf{*} \\
CN en moyenne de puissance 4 & +0,41\,\textbf{*} & +0,10\,\textbf{*} & +0,31\,\textbf{*} & +2,58\,\textbf{*} \\
\multicolumn{5}{@{}l}{\emph{Autres idées}} \\
règle temps-score propre à chaque circuit & $-$0,19\,\textbf{*} & $-$0,10\,\textbf{*} & $-$0,16\,\textbf{*} & $-$2,37\,\textbf{*} \\
score qui pénalise moins les grosses pertes & +0,04\,\textbf{*} & $-$0,04\,\textbf{*} & +0,19\,\textbf{*} & $-$0,44\,\textbf{*} \\
niveau estimé par maximum de vraisemblance & $-$0,98\,\textbf{*} & $-$0,50\,\textbf{*} & $-$1,27\,\textbf{*} & $-$5,69\,\textbf{*} \\
\multicolumn{5}{@{}l}{\textsc{Sprint}} \\
\multicolumn{5}{@{}l}{\emph{Références}} \\
CN officiel de la FFCO & $-$2,36\,\textbf{*} & $-$1,61\,\textbf{*} & $-$2,62\,\textbf{*} & $-$9,16\,\textbf{*} \\
Top linéaire : règle FFCO pour le circuit, CN en moyenne ordinaire & $-$0,05\,\textbf{*} & +0,05 & $-$0,12\,\textbf{*} & $-$2,80\,\textbf{*} \\
\multicolumn{5}{@{}l}{\emph{Exposant sur le temps, règle FFCO pour le circuit}} \\
exposant 1,5 & $-$0,17\,\textbf{*} & $-$0,13 & $-$0,11\,\textbf{*} & $-$1,18\,\textbf{*} \\
exposant 2 & $-$0,29\,\textbf{*} & $-$0,33\,\textbf{*} & $-$0,04 & $-$0,83\,\textbf{*} \\
exposant 3 & $-$0,61\,\textbf{*} & $-$0,72\,\textbf{*} & +0,01 & $-$1,29\,\textbf{*} \\
exposant 2 pour le circuit seulement & $-$0,30\,\textbf{*} & $-$0,29\,\textbf{*} & $-$0,18\,\textbf{*} & $-$3,54\,\textbf{*} \\
\multicolumn{5}{@{}l}{\emph{Règle FFCO pour le circuit, autre moyenne pour le CN}} \\
CN en moyenne cubique & +0,06\,\textbf{*} & +0,06\,\textbf{*} & +0,09\,\textbf{*} & +0,56\,\textbf{*} \\
\multicolumn{5}{@{}l}{\emph{Autres mesures du circuit}} \\
géométrique, deux tiers, CN en moyenne ordinaire & $-$0,01 & +0,10 & $-$0,15\,\textbf{*} & $-$2,73\,\textbf{*} \\
géométrique, deux tiers, CN en moyenne quadratique & +0,05\,\textbf{*} & +0,15\,\textbf{*} & +0,01 & +0,13\,\textbf{*} \\
arithmétique, tout le peloton, CN en moyenne ordinaire & $-$0,35\,\textbf{*} & $-$0,49\,\textbf{*} & $-$0,79\,\textbf{*} & $-$3,36\,\textbf{*} \\
arithmétique, tout le peloton, CN en moyenne quadratique & $-$0,33\,\textbf{*} & $-$0,51\,\textbf{*} & $-$0,74\,\textbf{*} & $-$0,64\,\textbf{*} \\
géométrique, moitié la plus rapide, CN en moyenne cubique & $-$0,02 & +0,02 & +0,24\,\textbf{*} & +0,25\,\textbf{*} \\
géométrique atténuant les temps aberrants (Huber), quadratique & +0,12\,\textbf{*} & +0,13 & $-$0,21\,\textbf{*} & +0,84\,\textbf{*} \\
\multicolumn{5}{@{}l}{\emph{Moyenne géométrique sur tout le peloton}} \\
géométrique, deux tiers, CN en moyenne cubique (retenu en sprint) & +0,12\,\textbf{*} & +0,15\,\textbf{*} & +0,11\,\textbf{*} & +0,72\,\textbf{*} \\
\multicolumn{5}{@{}l}{\emph{Autres idées}} \\
règle temps-score propre à chaque circuit & $-$0,14\,\textbf{*} & $-$0,13 & $-$0,15\,\textbf{*} & $-$3,41\,\textbf{*} \\
score qui pénalise moins les grosses pertes & $-$0,28\,\textbf{*} & $-$0,32\,\textbf{*} & $-$0,28\,\textbf{*} & $-$4,27\,\textbf{*} \\
niveau estimé par maximum de vraisemblance & +0,09\,\textbf{*} & +0,16 & +0,15\,\textbf{*} & +0,24\,\textbf{*} \\
\bottomrule \end{longtable}

\begin{table}[h]\centering\small
\caption{Le recalage en forêt : la moyenne des 20 \% meilleurs CN au 1\textsuperscript{er} octobre 2026 (niveau habituel, CN Top), pour la méthode d'alors et pour la méthode retenue ; cible 5\,600.}\label{tab:recal}
\begin{tabular}{@{}p{7.2cm}rrrr@{}}\toprule & \multicolumn{2}{c}{Méthode d'alors} & \multicolumn{2}{c}{Méthode retenue} \\ Recalage & Niveau & CN Top & Niveau & CN Top \\ \midrule
Jusqu'au 13 octobre : la moyenne des 13 mois qui se terminent 2 mois plus tôt & 6007 & 6144 & 6592 & 6757 \\
Une droite tracée sur ces 13 mois, lue le jour du facteur & 5816 & 5949 & 6022 & 6174 \\
La même droite, lue 6 mois plus loin & 5692 & 5821 & 5667 & 5812 \\
La moyenne des 13 mois, prolongée de 7 mois au rythme de la montée des 3 dernières années & 5716 & 5846 & 5723 & 5866 \\
La moyenne des 13 mois, prolongée de 7 mois au rythme de la montée des 10 dernières années (retenue) & 5663 & 5792 & 5639 & 5780 \\
\bottomrule\end{tabular}\end{table}

\begin{table}[h]\centering\small
\caption{Jeunes (20 ans et moins) et vétérans (55 ans et plus) face aux H/D\,21 : correction du CN qui mettrait le groupe à sa place (+ : le groupe est sous-évalué), 2012--2026, intervalle à 95 \%.}\label{tab:calib}
\begin{tabular}{@{}llrrr@{}}\toprule & & Officiel & Top quadratique & Méthode retenue \\ \midrule
Forêt & jeunes & +8,9 \% [+7,7 ; +10,1] & +3,3 \% [+2,7 ; +4,1] & +0,2 \% [$-$0,4 ; +0,8] \\
Forêt & vétérans & $-$7,8 \% [$-$8,7 ; $-$6,9] & $-$8,4 \% [$-$9,1 ; $-$7,7] & $-$9,1 \% [$-$9,8 ; $-$8,4] \\
Sprint & jeunes & +5,1 \% [+3,9 ; +6,2] & +2,3 \% [+1,5 ; +3,2] & +2,3 \% [+1,6 ; +3,2] \\
Sprint & vétérans & $-$17,2 \% [$-$18,2 ; $-$16,3] & $-$5,8 \% [$-$6,5 ; $-$5,0] & $-$5,6 \% [$-$6,3 ; $-$4,8] \\
\bottomrule\end{tabular}\end{table}
